\documentclass[10pt,a4paper]{amsart}
\usepackage[margin=19mm]{geometry}
\usepackage{amsmath,amssymb,mathtools}
\usepackage[scr=rsfs,cal=euler]{mathalfa}
\usepackage{xfrac}
\usepackage[unicode,hidelinks,bookmarks=false]{hyperref}
\newcommand{\ie}{i.e.\ }
\newcommand{\cf}{cf.\ }
\newcommand{\resp}{respectively\ }
\newcommand{\Subsection}{\subsection}
\providecommand{\nobreakdash}{\nobreak\hbox{-}}
\setlength{\emergencystretch}{2em}
\multlinegap0pt
\usepackage{amssymb}
\usepackage{tikz}
\usepackage{mathtools}
\newtheorem{theorem}{Theorem}
\newtheorem{lemma}{Lemma}
\DeclareMathOperator{\BHO}{BHO}
\newcommand{\avg}[1]{\left\langle {#1}\right\rangle}
\title{Weak-type estimates for the Bergman projection on planar domains}
\author{A. Walton Green, Cody B. Stockdale, Brandon Sweeting, Nathan A. Wagner}
\date{Source: arXiv:2607.09642v1 (CC BY 4.0)}
\begin{document}
\maketitle

\noindent\emph{Source and licence.} Weak-type estimates for the Bergman projection on planar domains. Version arXiv:2607.09642v1 (CC BY 4.0), distributed by arXiv under the Creative Commons Attribution 4.0 International licence, http://creativecommons.org/licenses/by/4.0/, as recorded in the arXiv licence metadata for this exact version and shown on the abstract page https://arxiv.org/abs/2607.09642v1. Full licence notice: \texttt{licenses/arxiv-2607.09642-CC-BY-4.0.txt}.\par\smallskip

\noindent\textit{Editorial scope.} Counted unit: the article's theorem thm:max-1 (source0.tex lines 498--499). The article's complete original proof of it is delivered with it (source0.tex lines 515--577, one \texttt{proof} environment of 63 lines, which the article places away from the statement). No mathematics is written, repaired or re-derived: the statement and the proof are the article's own lines, in the article's own notation, and no retained line is substituted or re-flowed.  Retained uncounted as the source-local context the proof needs: lines 501--507.  Nothing was omitted from any retained span, so the package carries no omission record.  No retained line contains a citation, so the wrapper carries no bibliography and no bibliography entry.  No retained line contains a figure environment, an image-inclusion command or drawing code, so no image is delivered and the package declares no asset dependency.  Editorial formatting only, in addition: an AMS \texttt{amsart} preamble that restates the article's own declarations and macros used by the retained lines, namely the environments \texttt{theorem}, \texttt{lemma} on separate counters, together with the standard packages those lines need.  The retained environments print in this document as Theorem 1, Lemma 1 in order of appearance, whereas the article prints the same items with the numbering of its own full document.  The complete native source of the article as distributed in the arXiv e-print for this exact version is bundled unchanged as \texttt{sources/204886/source0.tex}.
\begin{lemma}\label{lem:b1}
    If $u \in B_1\cap \BHO$, then there exists $\delta>0$ such that 
    \[
        |\{ z \in Q\colon u(z) > \lambda \}| \lesssim |Q|\left( \frac{\avg{u}_Q}{\lambda}\right)^{1 + \delta}
    \]
    for any Carleson cube $Q$ and $\lambda > 0$.
\end{lemma}

\begin{theorem}\label{thm:max-1} If $u$ belongs to $B_1 \cap \BHO$, then $M^u$ is bounded from $L^1(\mathbb{D}, u^2)$ to $L^{1,\infty}(\mathbb{D},u^2)$.
\end{theorem}

\begin{proof}[Proof of Theorem \ref{thm:max-1}]
    Let $u \in B_1$ and $f \in L^1(\mathbb{D},u^2)$. Recall the desired estimate
    \[
        \sup_{\lambda > 0} u^2(\{z \in \mathbb{D} \colon M^uf(z) > \lambda \}) \lesssim \frac{1}{\lambda}\|f\|_{L^1(\mathbb{D},u^2)}
    \]
    By homogeneity and the change of variables $f \mapsto uf$, it suffices to prove
    \[
        u^2(\{z \in \mathbb{D} \colon Mf(z) > u(z)\}) \lesssim \|f\|_{L^1(\mathbb{D},u)}.
    \]
    Moreover, by homogeneity, we may also assume that $u(\mathbb{D}) = 1$, and, by standard limiting arguments, that $f \in L^\infty(\mathbb{D})$ and $f \ge 0$. Let $a > 4$ and $q \in \mathbb{N}$ satisfying $a^{q-1} \le [u]_{B_1} < a^q$ (recall that $[u]_{B_1} \ge 1$). Therefore, as $u \in B_1$,
    \[
        1 = \int_\mathbb{D}u\,dA \le [u]_{B_1}\operatorname{ess\,inf}_{z \in \mathbb{D}} u(z) < a^q\operatorname{ess\,inf}_{z \in \mathbb{D}}u(z).
    \]
    It follows that $u(z) \ge a^{-q}$ a.e. on $\mathbb{D}$. Recall that for $f \in L^\infty(\mathbb{D})$ and any $k \in \mathbb Z$,
    \[
        \{z \in \mathbb{D} \colon Mf(z) > a^k \} = \bigcup_{j}Q_{k,j}
    \]
    for some collection of maximal dyadic Carleson boxes $\{Q_{k,j}\}_j$ for which $\avg{f}_{Q_{j,k}} \sim a^k$. Therefore, 
    \begin{align}
        \notag u^2(\{z \in \mathbb{D} \colon Mf(z) > u(z) \}) 
        &= \sum_{k = -q}^\infty u^2\left(\{ a^k < u \le a^{k+1}\} \cap \{ Mf > u \}\right)\\
        \notag &\le \sum_{k = -q}^\infty u^2\left(\{ a^k < u \le a^{k+1}\} \cap \{ Mf > a^k \}\right)\\
        &= \sum_{k = -q}^\infty \sum_{j}u^2(\{Q_{k,j} \colon a^k < u \le a^{k+1}\}). \label{e:B1-step-1}
    \end{align}
   These cubes are nested: for $-q \le m \le k$, each $Q_{k,j}$ is contained in a unique $Q_{m,i}$. We take the sum above over those cubes $Q_{k,j}$ for which $\{Q_{k,j} \colon a^k < u \le a^{k+1}\} \neq \emptyset$. For any such cube $Q_{k,j}$, we have
    \begin{equation}\label{e:avg_est}
        \avg{u}_{Q_{j,k}} \le [u]_{B_1} a^{k+1}
    \end{equation}
    We will sum the measures of these sets based on the $v$-averages on their associated cubes. We say a cube $Q_{k,j}$ is principal if $k = -q$, or if for the smallest principal cube containing $Q_{k,j}$, say $Q_{m,i}$, we have
    \[
        \avg{u}_{Q_{k,j}} \ge 2 \avg{u}_{Q_{m,i}}.
    \]
    Let $P := \{(k,j) \colon Q_{k,j} \text{ is principal}\}$ and, for $(k,j) \in P$ and $m \ge k$, let $F(k,j,m)$ denote the collection of indices $(m,i)$ for which $Q_{k,j}$ is the smallest principal cube containing $Q_{m,i}$. By the aforementioned nestedness, each cube is contained in some principal cube. Therefore, applying Lemma~\ref{lem:b1}, if $(m,i) \in F(k,j,m)$, then
    \[ u^2(\{Q_{m,i} \colon a^m < u \le a^{m+1}\})\lesssim a^{2m}|\{Q_{m,i} \colon a^m < u \le a^{m+1}\}| \] 
    \[ \lesssim a^{(1-\delta)m}|Q_{m,i}|\avg{u}_{Q_{m,i}}^{1+\delta} \lesssim a^{(1-\delta)m}|Q_{m,i}|\avg{u}_{Q_{k,j}}^{1+\delta} .
    \]
    By maximality of $Q_{k,j}$, we have that $Mf(z) = M(f\chi_{Q_{k,j}})(z)$ for each $z \in Q_{k,j}$. Consequently, by the weak-type $(1,1)$ estimate for the maximal operator, we have 
    \[
        \sum_{(m,i) \in F(k,j,m)}|Q_{m,i}| \le |\{ Q_{k,j} \colon M(f\chi_{Q_{k,j}}) > a^m \}| \lesssim a^{-m}\int_{Q_{k,j}}f\,dA.
    \]
   Therefore, combining the above two displays,
    with \eqref{e:avg_est}, we have
    \begin{align*}
        \eqref{e:B1-step-1}
        &\lesssim \sum_{(k,j) \in P}\sum_{m=k}^\infty \sum_{(m,i) \in F(k,j,m)} a^{m(1-\delta)}\avg{u}_{Q_{j,k}}^{1 + \delta}|Q_{m,i}| \\
        &\lesssim \sum_{(k,j) \in P} \avg{u}_{Q_{j,k}}^{1 + \delta}\left(\int_{Q_{k,j}}f\,dA\right)\sum_{m=k}^\infty a^{-m\delta}\\
        &\lesssim \sum_{(k,j) \in P} \avg{u}_{Q_{j,k}} \left(\int_{Q_{k,j}}f\,dA\right)\\   
        &= \int_{\mathbb{D}}f(z)\left(\sum_{(k,j) \in P}\avg{u}_{Q_{k,j}}\chi_{Q_{k,j}}(z)\right)\,dA(z).
    \end{align*}
    Since dyadic Carleson boxes have finite overlap, each $z \in \mathbb{D}$ belongs to at most finitely many such boxes, hence to at most finitely many principal cubes. For each $z \in \mathbb{D}$, let $Q_k(z)$ denote the principal cube at level $k$ containing $z$ (when one exists) and let $Q(z) = Q_{k_0}(z)$ denote the smallest such principal cube. Then
    \begin{align*}
        \sum_{(k,j) \in P}{|Q_{k,j}|}\avg{u}_{Q_{k,j}}\chi_{Q_{k,j}}(z) 
        &= \sum_{k = -q}^{k_0}\avg{u}_{Q_k(z)}\\
        &\le \avg{u}_{Q(z)}\sum_{k = -q}^{k_0}2^{-k}\\
        &\lesssim 2^q\avg{u}_{Q(z)} \\
        &\lesssim [u]_{B_1}^2u(z).
    \end{align*}
    Applying this estimate, we conclude
    \[
        u^2(\{z \in \mathbb{D} \colon Mf(z) > u(z) \}) \lesssim \|f\|_{L^1(\mathbb{D},u)},
    \]
    as desired.
\end{proof}

\end{document}
